\section{System Overview}
\label{sec:system}

The codebase is organized around four connected layers: syntax, static
well-formedness, executable semantics, and refinement. Together they support
the workflow ``write a small \llvm-like definition, execute it, and prove that
one definition is refined by another.''

\subsection{Abstract Syntax and Embedded Surface Syntax}

The core syntax lives in the \texttt{AST} modules. The project defines \lean
data types for widths, identifiers, types, expressions, instructions,
terminators, blocks, and single-block function definitions. The supported
instruction forms are intentionally small: integer binary operations,
selected integer conversions, and \texttt{freeze}. This is enough to express
the example optimizations and to surface some of the semantic issues that
make \llvm optimization correctness subtle.

One design decision that paid off early was to support a concrete surface
syntax inside \lean rather than constructing every example through raw
constructors. Quasiquoters are introduced to let the theorems look very
close to the \llvm syntax they are about, such as \texttt{[llvm-definition| ... ]}.
The syntax layer also supports symbolic widths,
written for example as \texttt{i\$0}, where the \texttt{\$0} is a
metavariable slot.
This means that a theorem can be stated
once and later instantiated at any concrete width.

The syntax support is not limited to parsing. The project also includes
printing and unexpansion support, so the embedded syntax is reasonably
round-trippable and readable in the \lean environment. This matters for a
project of this size: readability directly affects how easy it is to inspect
the generated terms and maintain the proof scripts.

\subsection{Symbolic Widths}

A central technical choice was to reuse the
\texttt{ConcreteOrMVar} datatype from \leanmlir to represent widths.
A width is either concrete, such as
\texttt{i8}, or a metavariable slot such as \texttt{i\$0}. Widths can then be
instantiated by a vector of natural numbers. This design gives the project a
useful middle ground. It is not fully dependent over width constraints, but
it is much more expressive than hard-coding one width at a time.

This choice directly supports one of the most attractive parts of the
project: width-generic theorems. Instead of proving the same identity again
for \texttt{i8}, \texttt{i16}, and \texttt{i32}, we can state one theorem over
a symbolic-width definition and instantiate it. This kept the project small
without making it trivial.

\subsection{Well-Formedness}

The semantics is only intended to run on structurally sensible definitions,
so we added explicit well-formedness predicates. Expressions are checked for
type compatibility and local scoping. Instructions are checked against the
expected operand types. Code checks enforce that instruction identifiers are
defined once and then become available in scope for later instructions.
Definitions check prototype shape, argument counts, duplicate argument names,
instruction well-formedness, and the compatibility of the terminator with the
declared return type.

This static layer is intentionally modest. It is not a full typing judgment
with theorem-proved preservation properties. However, it serves two practical
purposes. First, it makes the refinement theorem state cleaner by separating
semantic behavior from obviously malformed programs. Second, it prevents many
proofs from devolving into low-level bookkeeping about ill-scoped or
ill-typed examples.

\subsection{Executable Semantics}

The semantic layer interprets definitions in a state monad. The visible state
currently consists only of registers, stored in a hash map from identifiers to
register values. Register values are either bitvectors paired with a width or
\texttt{void}.

Expressions evaluate to integer values of a chosen width. Integer literals and
booleans are embedded into bitvectors, identifiers look up register state, and
\poison propagates through the integer semantics already provided in the
utility layer. Executing a definition binds its arguments, interprets the code
line by line, and then interprets the return terminator.

The semantics is executable in the strong sense that a fully instantiated
definition can be run from \lean and returns a concrete \texttt{Except}
result. This turned out to be valuable well beyond testing. Because the
semantic functions reduce under simplification, many small proof goals can be
discharged by computation rather than by manually replaying every semantic
step.

\subsection{Explicit \texttt{undef} Supply}

The newest and most interesting semantic feature is explicit \vundef
handling. Initially, \vundef was not supported in the executable semantics.
The final system introduces a separate supply state that stores a list of
natural numbers together with a current index. Encountering \vundef during
evaluation draws the next value from that supply and coerces it to the
requested bitvector width.

This is not intended to be a perfect model of \llvm's nondeterminism.
Instead, it is a pragmatic executable approximation that makes it possible to
state concrete refinement theorems involving \vundef and to inspect how many
\vundef choices a computation consumes. Importantly, this design forced the
refinement relation to become more explicit as well, because observable
behavior now depends on the chosen \vundef supply.

\subsection{Refinement}

The project defines refinement semantically on function definitions. In
essence, \texttt{x~$\sqsubseteq$~y} has two parts. First, if the source
definition \texttt{x} is defined for every \vundef supply, then the target
definition \texttt{y} must also be defined for every \vundef supply. Second,
every concrete return value produced by the target definition \texttt{y} must
either already be returned exactly by the source definition \texttt{x}, or be
covered by a poison-aware source return, assuming both definitions are well
formed, have compatible signatures, and are run on well-formed arguments.

For the current \vundef model, the definition quantifies over the \vundef
supplies used to run the two sides separately. This makes the refinement
relation slightly unusual, but it matches the explicit-supply execution model
and was sufficient to prove the theorems in the repository.

\small
\begin{leanlisting}
def Definition.IsRefinedBy (x y : @AST.Definition φ) : Prop :=
  ∀ (args : List RegisterValue),
    AST.Definition.WellFormed x →
    AST.Definition.WellFormed y →
    Definition.SignatureCompatible x y →
    Semantics.Definition.ArgValuesWellFormed x args →
    (((∀ (u : UndefChain), ∃ (ret_x : RegisterValue),
        runNanoLlvmStateM (evalNanoLlvmDefinition x args) default ⟨u, 0⟩ = .ok ret_x) →
      (∀ (u' : UndefChain), ∃ (ret_y : RegisterValue),
        runNanoLlvmStateM (evalNanoLlvmDefinition y args) default ⟨u', 0⟩ = .ok ret_y)) ∧
      (∀ (u : UndefChain) (ret : RegisterValue),
        runNanoLlvmStateM (evalNanoLlvmDefinition y args) default ⟨u, 0⟩ = .ok ret →
        ∃ u',
          runNanoLlvmStateM (evalNanoLlvmDefinition x args) default ⟨u', 0⟩ = .ok ret ∨
            ∃ ret_x,
              runNanoLlvmStateM (evalNanoLlvmDefinition x args) default ⟨u', 0⟩ = .ok ret_x ∧
                RegisterValue.IsRefinedBy ret_x ret))
\end{leanlisting}
\normalsize

This definition is worth reading from top to bottom. The outer quantification
ranges over concrete argument lists, together with the side conditions that
both definitions are structurally well formed, have compatible signatures, and
are run on well-formed source arguments. The body then splits into two
semantic clauses.

The first clause is the definedness condition. If the source definition
\texttt{x} can return some value for every \vundef supply, then the target
definition \texttt{y} must also be able to return some value for every
\vundef supply. The second clause is the value condition. It quantifies over
successful target executions and asks that every concrete return value
produced by \texttt{y} must either be reproduced exactly by \texttt{x},
possibly under a different \vundef supply, or be justified by a source return
that refines it through the poison-aware relation on register values. This
asymmetry is what allows the target to be more defined than the source
without introducing any genuinely new observable behavior.
